\section{Preliminaries about the anti-de Sitter space}
\label{Preliminary-AdS}

In this section, we collect some preliminary results about $\mathrm{AdS}^{3}$. We~refer to \cite[Section~9]{KaKob16} where they illustrate their general theory
for reductive symmetric spaces $X=G/H$ in details in the special setting where
$X=\mathrm{AdS}^{3}$. See also~\cite{Kan19}.

Let $Q$ be a quadratic form on $\mathbb{R}^{4}$ defined by
$Q(x) = x_1^2 + x_2^2-x_3^2-x_4^2$ for $x= (x_1,x_2,x_3,x_4)$ and
we set
\begin{gather*}
\mathbb{H}^{2,1}:=\big\{x=(x_1,x_2,x_3,x_4) \in \mathbb{R}^4 \mid Q(x) = 1\big\}
\cong \mathrm{SO}_{0}(2,2)/\mathrm{SO}_{0}(2,1).
\end{gather*}
The tangent space $T_{x}(\mathbb{H}^{2,1})$ at $x\in\mathbb{H}^{2,1}$
is isomorphic to the orthogonal complement $(\mathbb{R} x)^{\bot}$ with respect to $Q$.
Then $-Q|_{(\mathbb{R} x)^{\bot}}$ is
a quadratic form of signature $(2,1)$ on
$T_{x}(\mathbb{H}^{2,1}) \cong (\mathbb{R} x)^{\bot}$ and
thus $-Q$ induces a Lorentzian structure on $\mathbb{H}^{2,1}$
with constant sectional curvature $-1$.
The~$3$-di\-men\-sional anti-de Sitter space
\begin{gather*}
\mathrm{AdS}^{3}:=\mathbb{H}^{2,1}/\{\pm 1\}
\cong \mathrm{SO}_{0}(2,2)/(\{\pm1\}\times\mathrm{SO}_{0}(2,1)),
\end{gather*}
inherits a Lorentzian structure through
the double covering $\pi\colon\mathbb{H}^{2,1}\rightarrow\mathrm{AdS}^{3}$.

\subsection{Some coordinates and ``pseudo-balls''}
%\label{coord}
In this subsection, we work with coordinates on $\mathbb{H}^{2,1}$ and
consider ``pseudo-balls'' in $\mathrm{AdS}^{3}$. We~identify
$\mathbb{H}^{2,1}$ with $\mathrm{SL}(2,\mathbb{R})$ using the isomorphism
\begin{align}
\begin{array}{@{}ccc}
\mathbb{H}^{2,1} & \xrightarrow{\cong} & \mathrm{SL}(2,\mathbb{R}), \\
x=(x_1,x_2,x_3,x_4) & \longmapsto & \begin{pmatrix}
x_1 + x_4 & -x_2 + x_3 \\
x_2 + x_3 & x_1 -x_4
\end{pmatrix}\!.\label{isom}
\end{array}
\end{align}
For $t\geq0$ and $\theta\in\mathbb{R}$, we use the notations
\begin{gather}
\label{k-a}
k(\theta) = \begin{pmatrix}
\cos\theta & -\sin\theta \\
\sin\theta & \cos\theta
\end{pmatrix}\!,\qquad
a(t) = \begin{pmatrix}
{\rm e}^t & 0 \\
0 & {\rm e}^{-t}
\end{pmatrix}\!.
\end{gather}
We embed $\mathbb{H}^{2,1}$ into $\mathbb{C}^{2}$ by
\begin{gather}
\label{embed}
x\mapsto (z_1,z_2) = \big(x_1 + \sqrt{-1}x_2, x_3 + \sqrt{-1}x_4\big).
\end{gather}
We note that $z_{1}\neq0$ if $x\in\mathbb{H}^{2,1}$.
Via the identification (\ref{isom}), we have
\begin{gather}
\label{polar}
(z_1,z_2) = \big((\cosh t){\rm e}^{\sqrt{-1}(\theta_1+\theta_2)},
(\sinh t){\rm e}^{\sqrt{-1}(\theta_1-\theta_2)}\big),
\end{gather}
if $x = k(\theta_1)a(t)k(\theta_2) \in \mathrm{SL}(2,\mathbb{R})$
(a ``polar coordinate''). In~particular, we have
\begin{gather*}
\cosh 2t = x_1^2 + x_2^2 + x_3^2 + x_4^2.
\end{gather*}
Next, we consider pseudo-balls on $\mathrm{AdS}^3$, as a special case of
Kassel--Kobayashi~\cite{KaKob16} for
reductive symmetric spaces.
\begin{Definition}
\label{def-pball}
For $x = (x_1,x_2,x_3,x_4)\in\mathbb{H}^{2,1}$, $\|x\|\in\mathbb{R}_{\geq0}$
is defined by
\begin{gather*}
\cosh\|x\|:=x_1^2 + x_2^2 + x_3^2 + x_4^2\qquad (= \cosh(2t)).
\end{gather*}
This function is invariant under $x\mapsto -x$, hence defines
a function on $\mathrm{AdS}^{3}$,
to be also denoted by $\|\cdot\|$ (a ``pseudo-distance'' from the origin).
The compact set
\begin{gather*}
B(R) := \big\{y\in\mathrm{AdS}^{3} \mid \|y\| \leq R\big\}
\end{gather*}
is called a pseudo-ball of radius $R$.
\end{Definition}

\subsection[Square-integrable eigenfunctions of the Laplacian on the anti-de Sitter space] {Square-integrable eigenfunctions of the Laplacian \\on the anti-de Sitter space}\label{preliminary-spherical}

In this subsection, we consider square-integrable eigenfunctions of
$\square_{\mathrm{AdS}^{3}}$ with eigenvalues $\lambda_{m}=4m(m-1)$. We~recall from \cite[Section~9]{KaKob16} the following decomposition of the open subset $\{Q > 0\}$ of
the flat pseudo-Riemannian manifold $\mathbb{R}^{2,2}=\big(\mathbb{R}^{4},Q({\rm d}x)\big)$:
\begin{align*}
\begin{array}{@{}ccc}
\{Q > 0\} & \xrightarrow{\cong} & \mathbb{R}_{>0} \times \mathbb{H}^{2,1}, \\[1ex]
x & \longmapsto & \big(\sqrt{Q(x)}, x/\sqrt{Q(x)}\big).
\end{array}
\end{align*}
Let $r=\sqrt{Q(x)}$. Then one has, see \cite[p.~215]{KaKob16},
\begin{align}
\label{decomposition-of-laplacian}
-r^2\square_{\mathbb{R}^{2,2}} = -\left(r\frac{\partial}{\partial r}\right)^2 - 2r\frac{\partial}{\partial r} + \square_{\mathbb{H}^{2,1}}.
\end{align}
Let $m$ be a positive integer and $k$ be a non-negative integer. In~the coordinates (\ref{embed}),
the homogeneous function $z_{1}^{-(k+2m)}z_{2}^{k}$ of degree $-2m$
 is harmonic with respect to $\square_{\mathbb{R}^{2,2}}$,
hence its restriction to the submanifold $\mathbb{H}^{2,1}$ is
an eigenfuction of $\square_{\mathbb{H}^{2,1}}$ with eigenvalue $\lambda_{m}=4m(m-1)$
by~the formula (\ref{decomposition-of-laplacian}).
Moreover, it is square-integrable with respect to the measure
$\sinh(2t){\rm d}\theta_{1}{\rm d}t{\rm d}\theta_{2}$ in~the polar coordinate (\ref{polar})
induced from the Lorentzian metric on $\mathbb{H}^{2,1}$,
as in the $k=0$ case \cite[Section~9]{KaKob16}.
This $L^{2}$-eigenfunction is invariant under $(z_{1},z_{2})\mapsto(-z_{1},-z_{2})$,
hence defines a~real analytic $L^{2}$-eigenfunction on $\mathrm{AdS}^{3}$
with eigenvalue $\lambda_{m}$, to be denoted by $\psi_{m,k}$.
The discrete spectrum
$\mathrm{Spec}_{d}(\square_{\mathrm{AdS}^{3}})$ coincides with
$\{\lambda_{m} \mid m \in \mathbb{N}\}$
and $L^{2}_{\lambda_{m}}\big(\mathrm{AdS}^{3}\big)$ is generated by
$\psi_{m,0}$ and its complex conjugate $\overline{\psi_{m,0}}$
as a representation of $\mathrm{SO}_{0}(2,2)$ (see \cite[Claim~9.12]{KaKob16}).
By~(\ref{polar}), we~have
\begin{align}
\label{polarFJ}
\psi_{m,k}(\pi(x))
= {\rm e}^{-2\sqrt{-1}(m\theta_1 + (m+k)\theta_2)}\tanh^{k}t\cosh^{-2m}t
\end{align}
for $x=k(\theta_1)\,a(t)\,k(\theta_2)\in\mathbb{H}^{2,1}$. We~refer to $\psi_{m,k}$ as a spherical function of type $(-m,m+k)$
in~accordance with the action of
$\mathrm{SO}(2)\times\mathrm{SO}(2)$.

\subsection{Convergence of generalized Poincar\'e series}\label{preliminary-Poincare}
In this subsection, we explain the fact about the discrete spectrum
of locally symmetric spaces by Kassel--Kobayashi~\cite{KaKob16}
in our $\mathrm{AdS}^{3}$ setting. We~use the following notation.
\begin{Notation}\quad
\begin{itemize}\itemsep=0pt
\item
Let $\grave{\ }G=\mathrm{PSL}(2,\mathbb{R})=
\mathrm{SL}(2,\mathbb{R})/\{\pm1\}$ and
$G=\grave{\ }G\times\grave{\ }G$.
\item
Let $\grave{\ }K=\mathrm{PSO}(2)=\mathrm{SO}(2)/\{\pm1\}$ and
$K = \grave{\ }K\times\grave{\ }K$.
\item
Let $E$ and $\grave{\ }E$ be respectively the identity elements of $G$ and $\grave{\ }G$.
\end{itemize}
\end{Notation}

\begin{Remark}
The double covering $\mathrm{SO}_{0}(2,2)\rightarrow G$ induces
an isomorphism $\mathrm{AdS}^{3}\cong G/\mathrm{diag} \grave{\ }G$ $(\cong \grave{\ }G)$.
From now on, we consider only discontinuous groups $\Gamma$ for $\mathrm{AdS}^{3}$
which are discrete subgroups of $G$.
This is enough for our purpose.
\end{Remark}

In order to study
$\operatorname{Spec}_{d}\big(\square_{\Gamma\backslash\mathrm{AdS}^3}\big)$,
Kassel--Kobayashi~\cite{KaKob16} considered the convergence and non-vanishing of
generalized Poincar\'e series
\begin{gather}
\label{Poincare-series}
\varphi^{\Gamma}(\Gamma x):=\sum_{\gamma \in \Gamma} \varphi\big(\gamma^{-1}x\big)
\end{gather}
for $K$-finite square-integrable eigenfunctions $\varphi$ of $\square_{\mathrm{AdS}^{3}}$.
For this, they used an analytic estimate of $\varphi$ and
a geometric estimate of the number of $\Gamma$-orbits
\begin{gather}
\label{orbit-count}
N_{\Gamma}(x,R) := \#\{\gamma \in \Gamma \mid \gamma x \in B(R)\}
\end{gather}
in the pseudo-ball $B(R)$ for $R>0$.
Since the $\Gamma$-action is properly discontinuous and $B(R)$ is compact,
we have $N_{\Gamma}(x,R)<\infty$.

The convergence of generalized Poincar\'e series is proved by~\cite{KaKob16} as follows.
For $g \in G$ and a~function $f$ on $\mathrm{AdS}^{3}$,
$\ell_{g}^{*}f$ is defined by $\ell_{g}^{*}f(x)=f\big(g^{-1}x\big)$.
\begin{Fact}[Kassel--Kobayashi~\cite{KaKob16}]\label{Poincare}
Let $\Gamma \subset G$ be a discontinuous group for $\mathrm{AdS}^{3}$
satisfying the exponential growth condition
\begin{gather}\label{C}
\exists A,a > 0,\qquad
\forall x\in \mathrm{AdS}^{3},\qquad
\forall R>0, \qquad
N_{\Gamma}(x,R) < A{\rm e}^{aR}.
\end{gather}
Then, for any $K$-finite eigenfunction $\varphi$ of $\square_{\mathrm{AdS}^{3}}$ with eigenvalue
$\lambda_{m}$
and any $g\in G$,
if $m > a$, then
$(\ell_{g}^{*}\varphi)^{\Gamma}$ $($see $\eqref{Poincare-series})$ is continuous and square-integrable on
$\Gamma\backslash\mathrm{AdS}^{3}$ and
an eigenfunction of~$\square_{\Gamma\backslash\mathrm{AdS}^{3}}$ with eigenvalue
$\lambda_{m}$.
\end{Fact}

\begin{Remark}\label{remPoincare}\quad
\begin{enumerate}\itemsep=0pt
\item[(1)]
Fact~\ref{Poincare} does not assert the non-vanishing of the series
$(\ell_{g}^{*}\varphi)^{\Gamma}$ which is more involved.
Kassel--Kobayashi~\cite{KaKob16} proved that
there exists $g\in G$ such that
$(\ell^{*}_{g}\psi_{m,0})^{\Gamma}\neq 0$
for sufficiently large $m \in \mathbb{N}$.
\item[(2)]
%\label{remcount}
By \cite[Lemma~4.6.4]{KaKob16}, if a discontinuous group $\Gamma$ is sharp
in the sense of \cite[Definition~4.2]{KaKob16},
then $\Gamma$ satisfies the exponential growth condition (\ref{C}).
Moreover, Kassel~\cite{Ka09} and Gu\'eriataud--Kassel~\cite{GuKa17}
proved that finitely generated discontinuous groups for $\mathrm{AdS}^{3}$
are always sharp (see Fact~\ref{proper} below).
\item[(3)]
There exist discontinuous groups
which do not satisfy the exponential growth condition~(\ref{C}).
Indeed, for any increasing function $f\colon\mathbb{R}\rightarrow \mathbb{R}_{>0}$
and any $x\in\mathrm{AdS}^{3}$,
we~con\-st\-ructed
a discontinuous group $\Gamma_{f,x}$ for $\mathrm{AdS}^{3}$
satisfying $N_{\Gamma_{f,x}}(x,R) > f(R)$ for sufficiently large $R>0$ in~\cite{Kan19}.
\end{enumerate}
\end{Remark}
